\documentclass[11pt,a4paper]{article}

\usepackage[margin=2.5cm]{geometry}
\usepackage[T1]{fontenc}
\usepackage{lmodern}
\usepackage{hyperref}
\usepackage{listings}
\usepackage{booktabs}
\usepackage{enumitem}
\usepackage{xcolor}

\hypersetup{colorlinks=true, linkcolor=blue!50!black, urlcolor=blue!50!black}

\lstset{
  basicstyle=\ttfamily\small,
  breaklines=true,
  frame=single,
  backgroundcolor=\color{gray!8},
  rulecolor=\color{gray!50},
  columns=fullflexible
}

\title{VTPSA Website: Technical Handover Notes\\
\large Part 1: Architecture and Initial Setup}
\author{}
\date{12 November 2026}

\begin{document}
\maketitle

\section{Overview}
The VTPSA website is a static website that publishes results logs, club details,
documents (constitution, rulebooks, minutes, reports) and photo galleries.
All content lives as files in a Git repository. There is no database and no
server-side code. The site is rebuilt and published automatically whenever
changes are pushed to the \texttt{main} branch.

\begin{itemize}[nosep]
  \item \textbf{Live site:} \url{https://vtpsa.github.io}
  \item \textbf{Repository:} \url{https://github.com/VTPSA/vtpsa.github.io}
\end{itemize}

\section{Design decisions and rationale}
\begin{description}[leftmargin=0pt, style=nextline]
  \item[Static site instead of a database-backed web application.]
  The content volume is small (about 20 documents, 20 photos and a few result
  tables per year) and changes rarely. A static site has no running costs, no
  security surface and no moving parts to maintain.

  \item[GitHub Pages hosting instead of Firebase.]
  Since February 2026, Cloud Storage for Firebase requires the pay-as-you-go
  Blaze plan, which needs a payment card and has no hard spending cap. The
  organisation pays by EFT and has no card. GitHub Pages is free for public
  repositories and needs no payment details.

  \item[Content stored in Git.]
  Every change to results, documents or photos is a commit with an author and
  date, which gives a complete history and makes handover to a new maintainer
  straightforward.

  \item[Astro as the site generator.]
  Astro builds plain HTML pages from components and content files, and can
  generate pages directly from data such as club files and results tables.
\end{description}

\section{Technology stack}
\begin{center}
\begin{tabular}{@{}ll@{}}
\toprule
\textbf{Component} & \textbf{Choice} \\
\midrule
Site generator & Astro (exact version in \texttt{package.json}) \\
Runtime & Node.js 22.23.4 or newer (the build pipeline uses Node 24) \\
Package manager & npm (\texttt{package-lock.json} must stay committed) \\
Hosting & GitHub Pages \\
Build and deploy & GitHub Actions \\
Version control client & GitHub Desktop \\
Editor & Visual Studio Code \\
\bottomrule
\end{tabular}
\end{center}

\section{Repository and configuration}
The repository is owned by the GitHub organisation \texttt{VTPSA}, not by an
individual, so ownership survives a change of maintainer. It is public and
named exactly \texttt{vtpsa.github.io}; this special name makes GitHub publish
it at the root address with no path suffix.

Key files:
\begin{itemize}[nosep]
  \item \texttt{astro.config.mjs}: sets \texttt{site} to
        \texttt{https://vtpsa.github.io}. No \texttt{base} value is needed
        because the site is served from the root.
  \item \texttt{package.json} and \texttt{package-lock.json}: dependencies and
        their exact versions.
  \item \texttt{src/pages/}: each file here becomes a page of the site.
  \item \texttt{public/}: files copied to the site unchanged (for example PDFs
        and logos).
  \item \texttt{.github/workflows/deploy.yml}: the build and deploy pipeline.
\end{itemize}

\section{Deployment pipeline}
The workflow in \texttt{.github/workflows/deploy.yml} runs on every push to
\texttt{main} and can also be started manually from the Actions tab of the
repository. It has two jobs:
\begin{enumerate}[nosep]
  \item \textbf{build}: checks out the code, then runs the official
        \texttt{withastro/action}, which installs dependencies, builds the site
        and uploads the result. The Node version is set explicitly:
\begin{lstlisting}
- name: Install, build, and upload your site
  uses: withastro/action@v3
  with:
    node-version: 24
\end{lstlisting}
  \item \textbf{deploy}: publishes the built site using
        \texttt{actions/deploy-pages}.
\end{enumerate}
In the repository settings, \textbf{Settings $\rightarrow$ Pages $\rightarrow$
Source} must be set to \textbf{GitHub Actions}. Publishing a change normally
takes one to two minutes after the push.

\section{Local development workflow}
The project is kept outside OneDrive (for example in
\texttt{C:\textbackslash Projects}) because the \texttt{node\_}\texttt{modules}
folder contains thousands of small files that sync tools handle badly. The
normal cycle is:
\begin{enumerate}[nosep]
  \item Open the project folder in Visual Studio Code.
  \item Run the site locally from the project folder with
        \texttt{npm run dev} and view it at \texttt{http://localhost:4321}.
  \item Edit files and check the result in the browser.
  \item Commit and push using GitHub Desktop. The site then rebuilds and
        publishes automatically.
\end{enumerate}

\section{Problems encountered and their fixes}
\begin{center}
\small
\begin{tabular}{@{}p{0.36\textwidth}p{0.56\textwidth}@{}}
\toprule
\textbf{Problem} & \textbf{Cause and fix} \\
\midrule
Astro setup wizard looped in Git Bash &
Interactive prompts do not work in Git Bash. Use PowerShell and pass the
answers as command-line options. \\
\addlinespace
\texttt{npm} blocked in PowerShell &
Script execution is disabled by default on Windows. Fixed with
\texttt{Set-ExecutionPolicy -Scope CurrentUser -ExecutionPolicy RemoteSigned}. \\
\addlinespace
Build failed: ``Node.js v20 is not supported'' &
The Astro action defaults to Node 20 but this Astro version needs Node 22.23.4
or newer. Fixed by setting \texttt{node-version: 24} in the workflow. \\
\addlinespace
``Site not found'' at the root address &
The repository was not named \texttt{vtpsa.github.io}. Renamed it and re-ran
the workflow. \\
\bottomrule
\end{tabular}
\end{center}

\section{Notes for the next maintainer}
\begin{itemize}[nosep]
  \item Keep the GitHub organisation owned by at least two people, and use an
        organisational contact email where possible.
  \item Do not delete \texttt{package-lock.json}; the build pipeline relies on it.
  \item If a build fails, open the repository's Actions tab, click the failed
        run, and read the log of the failing step.
  \item Version numbers in the workflow file (for example
        \texttt{withastro/action@v3}) may need updating over time; check the
        current Astro deployment guide.
\end{itemize}

\end{document}
